\documentclass{beamer}
%\documentclass[xcolor=dvipsnames]{beamer}
\usepackage[spanish]{babel}
\usepackage[utf8]{inputenc}
\usepackage{graphicx}

\newcommand{\beamer}{\textsc{beamer}}
\newtheorem{definicion}{Definición}
\newtheorem{ejemplo}{Ejemplo}

%%%%%%%%%%%%%%%%%%%%%%%%%%%%%%%%%%%%%%%%%%%%%%%%%%%%%%%%%%%%%%%%%%%%%%%%%%%%%%%
\title[Trabajo de Fin de Grado]{SyTW: Aplicación para la elaboración y desliegue de cuestionarios}
\author {
Juan José Labrador González
}

\institute[ULL]{Escuela Superior de Ingeniería y Tecnología \\
                Departamento de Ingeniería Informática y de Sistemas \\
                Universidad de La Laguna}
\date[24-07-2014]{24 de Julio de 2014}
%%%%%%%%%%%%%%%%%%%%%%%%%%%%%%%%%%%%%%%%%%%%%%%%%%%%%%%%%%%%%%%%%%%%%%%%%%%%%%%

%\usetheme{Berlin}
\usetheme{Madrid}

%%%%%%%%%%%%%%%%%%%%%%%%%%%%%%%%%%%%%%%%%%%%%%%%%%%%%%%%%%%%%%%%%%%%%%%%%%%%%%%
\definecolor{pantone254}{RGB}{122,59,122}
\definecolor{pantone3015}{RGB}{0,88,147}
\definecolor{pantone432}{RGB}{56,61,66}
\setbeamercolor*{palette primary}{use=structure,fg=white,bg=pantone254}
\setbeamercolor*{palette secondary}{use=structure,fg=white,bg=pantone3015}
\setbeamercolor*{palette tertiary}{use=structure,fg=white,bg=pantone432}
\setbeamercolor*{palette sidebar primary}{use=structure,fg=pantone254}
\setbeamercolor*{palette sidebar tertiary}{use=structure,fg=pantone3015}
\setbeamercolor*{block title}{bg=pantone3015,fg=white}
\setbeamercolor*{alerted text}{fg=pantone432}
\setbeamercolor*{item projected}{fg=pantone254}
\setbeamercolor*{section in toc shaded}{use=structure,fg=structure.fg}
\setbeamercolor*{section in toc}{fg=pantone3015}
\setbeamercolor*{subsection in toc shaded}{fg=pantone3015}
\setbeamercolor*{subsection in toc}{fg=pantone432}

%%%%%%%%%%%%%%%%%%%%%%%%%%%%%%%%%%%%%%%%%%%%%%%%%%%%%%%%%%%%%%%%%%%%%%%%%%%%%%%
\begin{document}
  
%++++++++++++++++++++++++++++++++++++++++++++++++++++++++++++++++++++++++++++++  
\begin{frame}

  \includegraphics[width=0.15\textwidth]{img/ullesc.eps}
  \hspace*{7.5cm}
  \includegraphics[width=0.16\textwidth]{img/etsii.eps}
  \titlepage

%  \begin{scriptsize}
%    \begin{center}
%     Escuela Técnica Superior de Ingeniería Informática\\
%     Universidad de La Laguna
%    \end{center}
%  \end{scriptsize}

\end{frame}
%++++++++++++++++++++++++++++++++++++++++++++++++++++++++++++++++++++++++++++++  

%++++++++++++++++++++++++++++++++++++++++++++++++++++++++++++++++++++++++++++++  
\begin{frame}
  \frametitle{Índice}  
  \tableofcontents
\end{frame}
%++++++++++++++++++++++++++++++++++++++++++++++++++++++++++++++++++++++++++++++  

\section{Introducción}
\begin{frame}

\frametitle{Introducción}
 
\end{frame}
%++++++++++++++++++++++++++++++++++++++++++++++++++++++++++++++++++++++++++++++  

\section{Metodología}
\begin{frame}

\frametitle{Metodología}

\end{frame}
%++++++++++++++++++++++++++++++++++++++++++++++++++++++++++++++++++++++++++++++  

\section{Resultados}
\begin{frame}

\frametitle{Resultados}

\end{frame}
%++++++++++++++++++++++++++++++++++++++++++++++++++++++++++++++++++++++++++++++  

\section{Conclusiones}
\begin{frame}

\frametitle{Conclusiones}
\end{frame}

% 
% \section{Contexto}
% 
% 
% %++++++++++++++++++++++++++++++++++++++++++++++++++++++++++++++++++++++++++++++  
% \begin{frame}
% 
% \frametitle{Contexto}
% %
% 
% \begin{enumerate}[<+->]
% \item 
% Tercer curso del Grado en Ingeniería Informática
% \item 6 créditos: 3 de Teoría 1,5 de problemas y 1,5 de prácticas de laboratorio
% \item 
% El número de alumnos matriculados supera los ciento sesenta
% \item 
% Se dividen en dos grupos para las tres horas de teoría y problemas (mañana y tarde)
% \item 
% Se dividen en siete grupos para la hora de prácticas de laboratorio
% \item 
% Cuatro profesores
% \item Dos a cargo de los
% respectivos grupos de teoría y el resto compartiendo horas de laboratorio
% \end{enumerate}
% 
% \end{frame}
% %++++++++++++++++++++++++++++++++++++++++++++++++++++++++++++++++++++++++++++++  
% 
% \section{Competencias Transversales: Trabajo en Equipo}
% 
% \begin{frame}
% \frametitle{Competencias}
% %\includegraphics[width=1.0\textwidth]{img/paradigm-shift-cartoon.eps}
% \begin{itemize}
% \item
% \pause
% T10. Capacidad de integrarse rápidamente y {\color{blue} trabajar} eficientemente y rápidamente {\color{blue} en equipo} \ldots % Capacidad de integrarse rápidamente y {\color{blue} trabajar eficientemente en equipo}s unidisciplinares y de colaborar en un entorno multidisciplinar.
% \pause
% \item T12. Capacidad de {\color{blue} relación interpersonal} %Capacidad de {\color{blue} relación interpersonal}.
% \pause
% \item T16. Capacidad de {\color{blue} planificación y organización} del trabajo personal %Capacidad de {\color{planificación y organización} del trabajo personal.
% \end{itemize}
% 
% \end{frame}
% 
% \section{Control de Versiones}
% 
% \begin{frame}
% \frametitle{Control de Versiones}
% \begin{center}
% \includegraphics[width=1.0\textwidth]{img/versioncontrol1.eps}
% \end{center}
% \end{frame}
% 
% %++++++++++++++++++++++++++++++++++++++++++++++++++++++++++++++++++++++++++++++  
% \begin{frame}
% \frametitle{Control de Versiones}
% \begin{center}
% {\it Se llama {\color{red} control de versiones} a la gestión de los diversos cambios que se realizan sobre los elementos de algún producto o una configuración del mismo. 
% 
% Una {\color{blue} versión, revisión o edición de un producto}, es el estado en el que se encuentra el mismo en un momento dado de su desarrollo o modificación. }
% \end{center}
% \end{frame}
% 
% 
% \begin{frame}
% \frametitle{Control de Versiones: La Máquina del Tiempo}
% \begin{center}
% \includegraphics[width=1.0\textwidth]{img/timemachine.eps}
% \end{center}
% \end{frame}
% 
% \begin{frame}
% \frametitle{Control de Versiones: Git}
% \begin{center}
% \includegraphics[width=0.5\textwidth]{img/quote-linus.eps}
% 
% \includegraphics[width=0.4\textwidth]{img/gitdefinition.eps}
% \end{center}
% 
% \end{frame}
% 
% \begin{frame}
% \frametitle{Ramas: Universos Paralelos}
% \begin{center}
% \includegraphics[width=1.0\textwidth]{img/version-control-branch.eps}
% \end{center}
% \end{frame}
% 
% \begin{frame}
% \frametitle{Mezclando Universos Paralelos}
% \begin{center}
% \includegraphics[width=1.0\textwidth]{img/merge.eps}
% \end{center}
% \end{frame}
% 
% \section{Colaboración y Control de Versiones}
% \begin{frame}
% \frametitle{Colaboración y Git}
% \begin{center}
% \includegraphics[width=0.7\textwidth]{img/collaborating1.eps}
% \end{center}
% {\it 
% Cada usuario modifica la copia local ({\color{red} working copy}) 
% y cuando el usuario decide
% {\color{blue} compartir los cambios} ({\color{red} push}) 
% el sistema automáticamente intenta {\color{blue}combinar} ({\color{red} merge})
% las diversas modificaciones. 
% }
% %------------------------------------------------------------------------------
% \end{frame}
% 
% \begin{frame}
% \frametitle{GitHub: Facebook para Programadores}
% \begin{columns}
%   \begin{column}{0.5\textwidth}
%   \includegraphics[width=1.0\textwidth]{img/worldcollaboration.eps}
%   \end{column}
%   \begin{column}{0.5\textwidth}
%   \begin{center}
%   \includegraphics[width=0.8\textwidth]{img/github.eps}
%   \end{center}
%   \end{column}
% \end{columns}
% 
% GitHub is home to {\color{red} over 13.1 million repositories}, making it the largest code host in the world.
% \end{frame}
% 
% \begin{frame}
% \frametitle{Incidencias/Issues}
% \begin{center}
% \includegraphics[width=1.0\textwidth]{img/githubissues.eps}
% \end{center}
% %------------------------------------------------------------------------------
% \end{frame}
% 
% \section{Conflictos y Resolución de Conflictos}
% 
% 
% \begin{frame}
% \frametitle{Conflictos}
% \begin{center}
% \includegraphics[width=0.8\textwidth]{img/conflict-resolver.eps}
% \end{center}
% 
% {\it 
% El principal problema es la posible
% aparición de {\color{blue} conflictos} que deban ser solucionados manualmente 
% que resultan de
% las posibles inconsistencias que surgen al modificar el mismo fichero
% por varias personas no coordinadas. 
% }
% \end{frame}
% %++++++++++++++++++++++++++++++++++++++++++++++++++++++++++++++++++++++++++++++  
% 
% \section{Programación en Pareja (Pair Programming)}
% \begin{frame}
% \frametitle{Programación en Parejas}
% \includegraphics[width=1.0\textwidth]{img/pair_programming_1.eps}
% \end{frame}
% %++++++++++++++++++++++++++++++++++++++++++++++++++++++++++++++++++++++++++++++  
% 
% 
% \section{Evaluación}
% %++++++++++++++++++++++++++++++++++++++++++++++++++++++++++++++++++++++++++++++  
% \begin{frame}
% \frametitle{Evaluación}
% 
% \begin{center}
%     \includegraphics[width=0.6\textwidth]{img/ruqlsrc.eps}
% 
%     \includegraphics[width=0.5\textwidth]{img/ruqlbrowser.eps}
% \end{center}
% \end{frame}
% 
% \section{Conclusión: Control de Versiones para Todo y para Tod@s}
% %++++++++++++++++++++++++++++++++++++++++++++++++++++++++++++++++++++++++++++++  
% \begin{frame}
% \frametitle{Control de Versiones para Todo y para Tod@s}
% 
% \begin{center}
%     \includegraphics[width=0.7\textwidth]{img/estacontribucion.eps}
% \end{center}
% \end{frame}
% 
\end{document}
